\documentclass[runningheads]{llncs}

\usepackage[T1]{fontenc}
\usepackage{newtxtext}
\usepackage[varvw]{newtxmath}
\usepackage{graphicx}
\usepackage{url}

\begin{document}

\title{Evaluating Single-Use Ticket Enforcement and Revenue
Reconciliation in a Tourism Ticketing Platform}

\titlerunning{Ticket Enforcement and Revenue Reconciliation}

% Replace the author names, affiliations, and email addresses.
\author{
Muhammad Rizqi Firdaus\inst{1} \and
Amanda Eka Faraghita\inst{1}
}

\authorrunning{M. R. Firdaus and A. E. Faraghita}

\institute{
Electronic Engineering Polytechnic Institute of Surabaya, Surabaya, Indonesia \\
\email{rizqifrds@it.student.pens.ac.id}\\\email{amandaef098@it.student.pens.ac.id}}

\maketitle

% Write the abstract after the experiments are complete.
\begin{abstract}
Tourism ticketing platforms require reliable ticket validation and
consistent revenue reporting. This study evaluates single-use ticket
enforcement and revenue-summary consistency in MegiLoka, a prototype
tourism ticketing platform. The evaluation comprises ticket acceptance
and rejection tests, concurrent validation requests, and reconciliation
of synthetic transactions against an independent reference calculation.
% Replace the following sentence with actual results.
The experimental results and their implications will be reported after
the tests are completed. The study is limited to the documented
prototype and experimental conditions.

\keywords{
Electronic ticketing \and
QR code \and
Ticket validation \and
Concurrent requests \and
Revenue reconciliation \and
Tourism information systems
}
\end{abstract}

\section{Introduction}

Tourism ticketing platforms connect several activities in a visitor
service process. Visitors select a destination and purchase tickets,
while the platform records a transaction and issues a ticket for
admission. At the entrance, an operator validates the ticket before
granting access. Transaction records can also support reports that
summarize ticket sales and revenue by destination and date. The
reliability of such reports depends on consistent rules for payment
status, ticket issuance, cancellations, and the period in which
revenue is recorded.

Two requirements are particularly relevant to this lifecycle. First,
a ticket intended for one admission should not be accepted more than
once, including when two validators submit requests at nearly the
same time. Second, reported revenue should agree with transactions
that meet the study's eligibility rules. These requirements are
evaluated here as prototype scenarios; no claim is made that repeated
ticket use or revenue discrepancies have occurred at a particular
tourism agency. A valid message authentication code can help detect
changes to ticket data, but the platform must also manage ticket and
payment states consistently to decide whether admission is permitted.

Previous e-ticketing research has investigated ticket authenticity,
replay protection, and double spending in different settings. For
example, research on NFC-based public transport ticketing describes
protection against replay attacks~\cite{ref_nfc_ticketing}, while
work on reusable electronic tickets discusses online detection of
double-spending attempts~\cite{ref_reusable_ticket}. Research on
dynamic QR ticketing also considers copied codes and validation
histories~\cite{ref_dynamic_qr}. These approaches provide useful
points of comparison, although their technologies and intended usage
rules differ from those of a single-use tourism ticket. The present
study focuses on concurrent validation requests alongside the
consistency of the platform's revenue summaries.

This paper evaluates MegiLoka, a prototype tourism ticketing platform,
with respect to ticket validation and revenue-summary consistency.
The evaluation consists of three experiments: acceptance and rejection
of tickets under predefined conditions, two concurrent validation
requests for each of 30 distinct tickets, and reconciliation of
100 synthetic transactions against an independently calculated
reference. The study distinguishes payload integrity from single-use
enforcement and treats HMAC as an existing mechanism used by the
prototype rather than a new algorithm. Its contribution is an
evaluation of the ticket lifecycle and internal revenue summaries
under documented test conditions; it does not establish comprehensive
security or demonstrate an increase in regional revenue.

\section{Related Work}

\subsection{QR-Based Ticket Authenticity and Replay Prevention}

Electronic tickets can use cryptographic mechanisms to detect
unauthorized changes to ticket data. However, validating a ticket
payload and preventing the reuse of an authentic ticket are distinct
requirements. Previous research has examined replay protection in
electronic ticketing and QR-based systems~\cite{ref_nfc_ticketing,
ref_dynamic_qr,ref_qr_security}.

% Compare the selected studies in detail after reading their complete
% methods and evaluation sections. Identify the payload validation
% mechanism, replay-prevention mechanism, and experimental conditions.

\subsection{Single-Use Redemption and Concurrent Validation}

A ticketing system must apply its usage rules consistently when
multiple validation requests are received. Online validation can
support the detection of repeated use, although the applicable rules
depend on the ticketing model~\cite{ref_reusable_ticket}.

% Explain how the selected studies handle ticket state, simultaneous
% requests, and atomic state changes. Do not claim that a study lacks
% a feature unless the complete source supports that conclusion.

\subsection{Transaction and Revenue Reconciliation}

Revenue reconciliation compares reported values with a reference
derived from eligible transactions. The comparison requires explicit
rules for payment status, cancellation, refunds, transaction dates,
and grouping by destination.

% Add verified primary literature or authoritative sources on
% transaction and payment reconciliation. Distinguish internal report
% reconciliation from reconciliation with bank or payment-provider
% settlement records.

\subsection{Positioning of This Study}

The NFC architecture addresses replay in a public-transport setting,
while the reusable-ticket study considers online double spending for
journeys that may include transfers. Dynamic QR research discusses
copied codes and validation histories. The RSA-PSS evaluation examines
tampering and replay in a digital-identity context. These studies
provide different points of comparison for the present prototype.

\begin{table}
\caption{Related work and relevance to this study.}
\label{tab:related-comparison}
\centering
\begin{tabular}{|p{2.4cm}|p{3.0cm}|p{4.0cm}|}
\hline
\textbf{Study} & \textbf{Focus} & \textbf{Boundary} \\
\hline
NFC ticketing~\cite{ref_nfc_ticketing} & Replay protection &
Different technology and transport setting. \\
\hline
Reusable ticket~\cite{ref_reusable_ticket} & Online double-spending
detection & Allows transfers under different usage rules. \\
\hline
Dynamic QR~\cite{ref_dynamic_qr} & QR copying and validation history &
Concurrent redemption and revenue evaluation are not established
from the reviewed source. \\
\hline
QR security~\cite{ref_qr_security} & Tampering and replay &
Digital identity rather than tourism revenue. \\
\hline
Present study & Validation and internal recap & Proposed tests use
ticket cases, concurrent requests, and synthetic transactions. \\
\hline
\end{tabular}
\end{table}

This study examines ticket payload integrity, single-use enforcement
under concurrent requests, and the consistency of revenue summaries
within one prototype. Its reconciliation experiment compares the
application's internal reports with an independently calculated
reference for synthetic transactions.

% Add a comparison table after completing the literature review.
% Suggested columns: study, payload validation, usage state,
% concurrent redemption, revenue reconciliation, and evaluation method.

\section{System and Research Methodology}

\subsection{Study Context and Requirements Evidence}

MegiLoka is a prototype tourism ticketing platform that includes
ticket booking, ticket validation, and revenue reporting. The present
study evaluates the prototype under controlled conditions.

% If observations, interviews, or procedural documents are used,
% identify the source, date, and relevance of each item. If no field
% evidence is available, retain the prototype-scenario framing and do
% not present the requirements as observed institutional problems.

\subsection{Ticket and Payment Lifecycle}

The lifecycle considered in this study comprises transaction
creation, payment-status recording, ticket issuance, ticket validation,
and revenue reporting. Payment status and ticket status are represented
separately. A transaction may have a payment status such as paid,
pending, or cancelled, while an individual ticket may have a status
such as valid, used, expired, or cancelled.

% Insert a lifecycle diagram here. Explain the actual transitions in
% the final version of the system, including cancellation and refund
% behavior when supported.
%
% \begin{figure}
% \centering
% \includegraphics[width=\textwidth]{ticket-lifecycle.pdf}
% \caption{Ticket and payment lifecycle in the evaluated platform.}
% \label{fig:lifecycle}
% \end{figure}

\subsection{QR Payload Integrity and Ticket Validation}

The prototype generates a QR payload containing a ticket code,
destination identifier, validity date, and an HMAC-based value. The
HMAC-SHA256 calculation covers the ticket code, destination identifier,
and validity date. The implementation stores the first 16 hexadecimal
characters of the resulting value in the payload.

Payload verification and ticket usage are separate checks. Payload
verification examines whether the signed fields match the expected
HMAC value. The database stores the ticket status used to determine
whether a ticket has already been redeemed.

% Document the exact validation logic of the tested version.
% The reviewed implementation also permits manual ticket-code input
% without a signed QR payload. Explain whether this path is in scope.
% Describe payment-status and validity-date checks only if they are
% present in the tested version.
%
% If an atomic update or another concurrency control mechanism is
% implemented, describe it here and identify the tested code version.
% Do not claim atomic enforcement before verifying the implementation.

\subsection{Revenue Definition and Independent Reference Calculation}

In the reviewed implementation, the validator verifies the payload,
looks up the ticket, rejects a used, expired, or cancelled status,
and otherwise marks the ticket used before committing. The status read
and update are separate operations; an atomic conditional update is
not apparent in the reviewed path. Concurrent single-use enforcement
therefore remains a question for testing.

The reviewed validator does not check transaction payment status or
compare the validity date with the current date. A manual ticket-code
format is accepted without a signed QR payload. These observations
come from code review and are not measured experimental results. The
tested version and any subsequent changes must be identified.

For this study, revenue must be defined using explicit transaction
eligibility and recording-date rules. The same date and destination
filters must be applied to the application report and to the
independent reference calculation.

% Specify:
% 1. Which payment statuses contribute to revenue.
% 2. How cancellations and refunds are treated.
% 3. Which date determines the reporting period.
% 4. Whether amounts are gross or net of refunds.
% 5. The currency and rounding rules.
%
% The independent reference must be calculated outside the
% application's recap function, using transaction-level data.

The number of transactions and the number of tickets are reported as
different measures because a transaction can contain more than one
ticket. The evaluated version of the platform stores ticket counts in
its recap; transaction counts therefore require a separate
transaction-level calculation or an explicitly implemented reporting
feature.

\subsection{Experimental Setup and Data Collection}

For the proposed baseline, only paid transactions not subsequently
cancelled contribute their recorded amount to the transaction date
and destination. Pending and cancelled transactions contribute zero.
The time zone, rounding, and cancellation rule must be frozen before
execution. Refunds are outside this baseline unless supported and
documented. For date $d$ and destination $g$, let $E_{d,g}$ be the
independently identified eligible transactions and $a_i$ their amounts.
The reference count and amount are
\begin{equation}
N^{\mathrm{ref}}_{d,g}=|E_{d,g}|,\qquad
R^{\mathrm{ref}}_{d,g}=\sum_{i\in E_{d,g}} a_i.
\end{equation}
The amount difference is
$\Delta_{d,g}=R^{\mathrm{sys}}_{d,g}-R^{\mathrm{ref}}_{d,g}$.
The reference is maintained separately from the application's recap.

The reviewed booking and point-of-sale paths create transactions as
paid and update a stored recap on the creation date. The report reads
that recap. Cancellation adjustments and refunds are not established
in those paths and require testing or a stated limitation.

The experiments are conducted using a documented version of the
prototype and a controlled test database. The test environment,
database engine, server configuration, data preparation procedure,
and logging mechanism must be recorded to support reproducibility.

% Record the application version or commit, operating system,
% database engine and version, server configuration, test dates,
% and time zone.
%
% Use a clean test database or document the baseline. Existing
% synthetic recap data must not be mixed with experimental results.
% State explicitly that payments are simulated.

\begin{table}
\caption{Experimental datasets and conditions.}
\label{tab:experimental-conditions}
\centering
\begin{tabular}{|p{4.0cm}|p{5.0cm}|}
\hline
\textbf{Experiment} & \textbf{Dataset and conditions} \\
\hline
Ticket acceptance and rejection &
At least 10 tickets per predefined category. \\
\hline
Concurrent validation &
30 distinct valid tickets and two validator sessions. \\
\hline
Revenue reconciliation &
100 synthetic transactions across defined statuses, dates,
destinations, and amounts. \\
\hline
\end{tabular}
\end{table}

\subsection{Evaluation Scenarios and Metrics}

\subsubsection{Experiment 1: Ticket Acceptance and Rejection.}

At least 10 tickets are prepared for each of the following categories:
valid, modified payload, already used, cancelled, and unpaid.
Expired tickets are added when the relevant rule is supported and
defined. The expected response and initial state are documented
before each test. For a modified payload, one signed field is changed
without generating a new HMAC value.

Each ticket is submitted for validation, and both the response and
the resulting usage state are recorded. Rejected tickets must not
increase the recorded number of admissions.

% Distinguish an explicit expired status from a validity date that
% has passed. Add separate categories when both are tested.

\subsubsection{Experiment 2: Concurrent Ticket Validation.}

Thirty distinct valid tickets are tested. For each ticket, two
validator sessions submit validation requests concurrently. Both
responses are recorded, followed by an inspection of the persistent
ticket state and available usage records.

The expected outcome for each round is one successful response, one
rejection, and one recorded use. The number of rounds with two
successful responses is divided by 30 to report the observed
double-acceptance rate.

% Record how concurrent requests are synchronized and measure latency
% for both sessions. A single used status alone cannot establish that
% only one request received a successful response. Preserve both
% responses and any available per-request audit records.

\subsubsection{Experiment 3: Transaction and Revenue Reconciliation.}

A dataset of 100 synthetic transactions is prepared with different
amounts and predefined payment statuses. The dataset includes paid,
pending, and cancelled transactions, as well as refunds if the
evaluated system supports them.

An independent reference is calculated by date and destination using
the study's revenue rules. Application results are obtained using the
same filters and compared with the reference for each group. Any
difference is investigated using transaction identifiers.

% Record the exact data-entry procedure. Explain whether transactions
% are created through the application or inserted for test setup and
% how the associated recap is generated.

\section{Results}

% Replace every placeholder in this section with observed results.
% If a planned experiment is changed or omitted, revise the methods,
% abstract, title, and claims accordingly.

\subsection{Ticket Acceptance and Rejection}

Table~\ref{tab:ticket-results} summarizes the observed responses for
each predefined ticket category.

\begin{table}
\caption{Ticket acceptance and rejection results.}
\label{tab:ticket-results}
\centering
\begin{tabular}{|p{2.5cm}|c|c|c|c|c|}
\hline
\textbf{Category} & \textbf{N} & \textbf{CA} & \textbf{CR}
& \textbf{IA} & \textbf{IR} \\
\hline
Valid & 10 & -- & -- & -- & -- \\
\hline
Modified payload & 10 & -- & -- & -- & -- \\
\hline
Already used & 10 & -- & -- & -- & -- \\
\hline
Cancelled & 10 & -- & -- & -- & -- \\
\hline
Unpaid & 10 & -- & -- & -- & -- \\
\hline
% Add expired categories if applicable.
\end{tabular}

\smallskip
\parbox{0.9\textwidth}{\footnotesize
N: number tested; CA: correctly accepted; CR: correctly rejected;
IA: incorrectly accepted; IR: incorrectly rejected.}
\end{table}

% Explain the observed results and identify any rejected ticket that
% changed the recorded usage state.

\subsection{Concurrent Ticket Validation}

Table~\ref{tab:concurrency-results} records the responses and
persistent usage observations for each concurrent-validation round.

\begin{table}
\caption{Concurrent validation results.}
\label{tab:concurrency-results}
\centering
\begin{tabular}{|c|c|c|c|c|c|}
\hline
\textbf{Round} & \textbf{Ticket ID} & \textbf{A} & \textbf{B}
& \textbf{Uses} & \textbf{Latency A/B} \\
\hline
1 & -- & -- & -- & -- & -- \\
\hline
% Add the remaining rounds or provide the complete table as
% supplementary material, subject to venue requirements.
\end{tabular}
\end{table}

% Report:
% - Rounds with two successful responses: [number] / 30.
% - Rounds with exactly one acceptance, one rejection, and one
%   recorded use: [number] / 30.
% - Latency summary and any incomplete or failed requests.

\subsection{Transaction and Revenue Reconciliation}

Table~\ref{tab:reconciliation-results} compares the system results
with the independent reference by date and destination.

\begin{table}
\caption{Transaction and revenue reconciliation by group.}
\label{tab:reconciliation-results}
\centering
\begin{tabular}{|p{1.5cm}|p{1.8cm}|c|c|c|c|c|}
\hline
\textbf{Date} & \textbf{Destination} & \textbf{RT} & \textbf{ST}
& \textbf{RA} & \textbf{SA} & \textbf{Diff.} \\
\hline
-- & -- & -- & -- & -- & -- & -- \\
\hline
% Add one row for each date and destination group.
\end{tabular}

\smallskip
\parbox{0.9\textwidth}{\footnotesize
RT: reference transaction count; ST: system transaction count;
RA: reference amount; SA: system amount; Diff.: system amount
minus reference amount.}
\end{table}

% Explain how system transaction counts were obtained. The existing
% recap's ticket count is not automatically a transaction count.
% Trace discrepancies to transaction IDs and report group-level
% results even when the grand total difference is zero.

\subsection{Failure Cases and Observed Causes}

Table~\ref{tab:failure-cases} presents representative cases in which
the observed behavior differs from the predefined expectation.

\begin{table}
\caption{Representative failure cases and observed causes.}
\label{tab:failure-cases}
\centering
\begin{tabular}{|p{1.5cm}|p{2.4cm}|p{2.4cm}|p{3.0cm}|}
\hline
\textbf{Case ID} & \textbf{Expected} & \textbf{Observed}
& \textbf{Cause supported by evidence} \\
\hline
-- & -- & -- & -- \\
\hline
% Add observed cases only. If no failures occur, state that none were
% observed under the tested conditions.
\end{tabular}
\end{table}

\section{Discussion}

\subsection{Ticket Authenticity and Replay Prevention}

A rejected modified payload addresses the tested integrity condition;
rejection of a second request addresses the observed usage sequence.
Two successful concurrent responses would indicate double acceptance
in that round even if one final ticket row is marked used. These tests
do not cover every attack path.

% Discuss payload integrity, sequential reuse, and concurrent
% validation separately. Explain how the observed responses relate to
% the implemented validation logic and database state changes.
% Do not infer that HMAC alone prevents reuse.

\subsection{Reliability of Revenue Summaries}

A zero grand-total difference can conceal discrepancies that cancel
across groups. Date-destination comparisons and transaction identifiers
are needed to identify such errors. Agreement against synthetic data
supports internal consistency only under the defined rules.

% Explain agreement or disagreement between the application and the
% independent reference. Discuss discrepancies per date and
% destination, including errors that cancel out in the grand total.

\subsection{Implications for Tourism Ticketing Services}

Traceable validation and reporting can support an operator's review
of admissions and transaction summaries. Claims about operational
benefits depend on measured outcomes and field evidence.

% Explain how the observed behavior affects ticket inspection,
% admission records, and the interpretation of revenue summaries.
% Relate the discussion to the prototype and evidence available.

\subsection{Limitations and Threats to Validity}

The experiments use synthetic transactions and simulated payment
states. Their results are limited to the tested application version,
dataset, environment, and operating rules. They do not establish
comprehensive system security, confirm settlement with a payment
provider, or demonstrate an increase in regional revenue.

% Add limitations concerning sample size, concurrent request
% synchronization, database engine, available audit logs, and
% field evidence as appropriate.

\section{Conclusion}

% Write this section after the experiments are complete.
%
% Answer the research question using actual values:
% 1. Ticket acceptance and rejection results by category.
% 2. Concurrent rounds with double acceptance out of 30.
% 3. Revenue and transaction differences by date and destination.
%
% State the implications and the limits of the evidence. Do not claim
% successful enforcement or accurate reconciliation if the results do
% not support those conclusions.

\begin{credits}

% Remove this item if no acknowledgment is required.
% \subsubsection{\ackname}
% The authors thank [name or institution] for [contribution].

\subsubsection{\discintname}
% Confirm this disclosure with all authors before submission.
The authors have no competing interests to declare that are relevant
to the content of this article.

% Replace this statement if there are relevant competing interests.

\end{credits}

\begin{thebibliography}{8}

\bibitem{ref_nfc_ticketing}
A Simple and Secure E-Ticketing System for Intelligent Public
Transportation Based on NFC.
\url{https://eudl.eu/doi/10.4108/icst.urb-iot.2014.257244}

\bibitem{ref_reusable_ticket}
A Reusable Public Transport Electronic Ticket System with Fast
Validation.
\url{https://doi.org/10.1007/s12469-024-00381-8}

\bibitem{ref_dynamic_qr}
Dynamic QR Codes for Ticketing Systems.
\url{https://fenix.tecnico.ulisboa.pt/downloadFile/844820067128428/87678-jose-oliveira-dissertacao.pdf}

\bibitem{ref_qr_security}
Slameto, A.A., Sutanto, Y., Kurniawan, H.:
Multi-Scenario Security Evaluation of an RSA-PSS Anti-Replay QR Code
System with a PRISMA 2020 Review of Indonesian Prior Work.
\url{https://doi.org/10.30871/jaic.v10i5.14095}

% Complete author names, publication details, page numbers, and
% access dates after verifying each source against the venue's
% bibliography requirements.

\end{thebibliography}

\end{document}
